\section{MARCO TEÓRICO}


En esta sección se citan los autores que han tenido influencia directa en tu investigación. Recuerda, debes escoger solo un método para realizar las citas y referencias, es decir, debes seleccionar entre gestores especializados como Mendeley, Zotero, EndNote, etc., Microsoft Word, o “Manuales”; no se deben mezclar entre sí. Nuestra recomendación principal es Mendeley. Evita referenciar sitios como blogs, Wikipedia, Rincón del Vago, Monografías.com y demás portales web que no se consideran fuentes primarias. No limites tu búsqueda a una sola herramienta (por ejemplo, solo www.google.com). Realiza búsquedas en diferentes plataformas académicas, tales como:

\begin{itemize}
    \item \textbf{Catálogo Sistema de Bibliotecas UdeA:} material impreso que reposa en Bibliotecas UdeA, tales como libros, revistas, tesis, diccionarios, informes, etc.
    \item\textbf{Bases de datos suscritas de la Biblioteca:} plataformas digitales con millones de documentos en texto completo.
    \item\textbf{Bases de datos de libre acceso:} Google Scholar, Microsoft Academic, Google Books, Redalyc, Scielo, Dialnet, DOAJ, PubMed, Base Search, entre muchas más.
    \item\textbf{Documentos con acceso restringido:} si requieres el texto completo de artículos o libros con acceso restringido, que por lo general se encuentran en bases de datos no suscritas por la Universidad de Antioquia, solicítalos en tu Biblioteca enviando título exacto, el DOI o la url del documento. Tenemos convenios nacionales e internacionales que nos permiten acceder a esta información.
\end{itemize}

Ejemplo de cita parafraseada, es decir, frase no textual adaptada con las palabras de quien escribe; esta forma de citación es la más adecuada en textos académicos, demuestran lectura, análisis y redacción propia \cite{sec:pediatria}. Ejemplo de “Cita textual menor a 40 palabras, al interior del párrafo; no utilice recurrentemente esta forma de citación, pues demuestra poco análisis y redacción propios” \cite[p. 9]{opc:ramirez}. Otros ejemplos aceptados en estilo IEEE 2020:

En citas paráfrasis, es posible mencionar o no el o los autores en la oración, ejemplos: resultados demostrados por Arango \cite{sec:pediatria}, resultados demostrados mediante publicación científica \cite{sec:pediatria}. Cita en paráfrasis de dos autores, con o sin autor mencionado: datos suministrados por Ramírez y Guzmán \cite{opc:ramirez}, datos suministrados en revista académica \cite{opc:ramirez}. Cita en paráfrasis con tres o más autores con et al.: resultados publicados por Baker et al. \cite{art:baker}, resultados publicados en revista científica \cite{art:baker}. Cita con dos fuentes, cada una en corchetes y en orden numérico: ambos autores coinciden en estos datos \cite{opc:bbva}, \cite{opc:film}. Cita con tres fuentes o más, con corchetes para el primero y el último, en orden numérico y separados por guion: diversos autores coinciden en estos datos \cite{opc:bbva}–\cite{lib:ieee}. 

Cita “textual menor a 40 palabras” \cite[p. 466]{art:espectador}, cita “textual menor a 40 palabras con páginas continuas o discontinuas” \cite[pp. 15-16]{tes:tesis}.  %[8, pp. 15-16). debe terminar en parentesis

Cita de cita en paráfrasis mencionando o no autores: Quintero y González, citados por Rioja \cite{con:conferencia}, Quintero y González, citados por \cite{con:conferencia}. Cita textual mayor a 40 palabras sin comillas:


    \begin{quotation}
    Es importante destacar que la revisión realizada permitió definir el constructo a evaluar, es decir, especificar el concepto de la e-inclusión que se asume en la investigación, así como los factores que deben ser considerados para su evaluación. Lo anterior constituye el fundamento conceptual de la tesis y es la base para desarrollarla \cite[p. 107]{art:acimed}. 
    \end{quotation}

Cita textual o directa con más de 40 palabras (se omiten las comillas), bloque aparte, sangría 1.25 cms. Ya que IEEE no señala nada en particular de este caso, se adaptan de las normas APA. Datos estadísticos analizados con SPSS \cite{opc:ibm}. Procure no incurrir en la citación excesiva:


\subsection*{MÉTODO PARA ELABORAR CITAS Y REFERENCIAS}
Recuerda, debes escoger solo un método para realizar las citas y las referencias, es decir, debes seleccionar entre un gestor especializado como Mendeley, Zotero o EndNote (confiabilidad alta), el gestor nativo Microsoft Word (confiabilidad media) o “manuales” (confiabilidad baja); no se deben mezclar entre sí, nuestra recomendación principal es Mendeley, pues te permite almacenar, subrayar e insertar notas en PDF´s, guardar fichas bibliográficas, elaborar citas y bibliografías en APA, Vancouver, IEEE y más de 1.000 estilos diferentes, exportar registros, compartir documentos, etc. Para mayor información de Mendeley.com y \LaTeX visitar:

\url{http://tiny.cc/zzrnuz}

La presente plantilla incluye los archivos requeridos para la construcción apropiada de las referencias bibliograficas bajo los parametros de las normas IEEE, todo se limita a modificar los campos del archivo \textit{refe.bib} segun las necesidades de su escrito. Aqui algunos ejemplos de los diferentes tipos de documentos referanciables.


Refencia a un libro \cite{lib:libro}, a un artículo \cite{art:baker}, a una seccion de libro sin título \cite{sec:pediatria}, a un acta \cite{act:acta}, a una conferencia \citep{con:conferencia}, a un capítulo o páginas específicas con título propio de un libro\citep[pp. 7-10]{ent:entre_libro},  a un documento sin editorial \citep{fll:folleto}, a las memorias de una conferencia \citep{ina:articulo_en_acta}, a un informe técnico \citep{inf:informe_tecnico}, a un manual \citep{man:manual}, a una tesis doctoral \citep{phd:phd_tesis}, a una tesis de maestría \citep{tes:tesis}, a una web o información de internet \citep{web:web} y para otros tipos de referencias diferentes a las anteriores \cite{opc:film}. Para el listado de referencias, tenga en cuenta que el apellido principal será la última palabra en el campo \textit{author}; Si como autor no figura una persona sino una organización, al llenar el campo \textit{author} del archivo \textit{.bib}, use $\{\}$ dobles (ver referencia \citep{lib:ieee}).

A continuación se muestra una de las referencias del archivo \textit{refe.bib}, donde \linebreak \textit{ent:entre\_libro} es el identificador de la misma y para referenciarla se debe escrbir el comando $\backslash cite\{ent:entre\_libro\}$ 

\newpage
\begin{verbatim}
@inbook{ent:entre_libro,
     author         = {Angela Zapata},
     title          = {Titulo de libro},
     chapter        = {},
     pages          = {10-20},
     publisher      = {Editor},
     volume         = {},
     series	        = {},
     type           = {},
     address        = {},
     edition        = {},
     year           = {20XX},
     month          = {},
     note           = {},
     language       = {},
     }
\end{verbatim}

Recuerde que \LaTeX , en la sección de referencias, solo muestra las que se han citado en su escrito; agregar la información en el archivo \textit{refe.bib} no implica que la referencia se muestre en su documento

Por otro lado, para citar ecuaciones e información semejante dentro del texto use $\backslash ref\{\}$, por ejemplo (\ref{eq:maxwell}) o (\ref{eq:maxwell3})

\begin{equation}\label{eq:maxwell}
    \oint_C \vec{B}\cdot d \vec{l}=\mu_0\int_S \vec{J}\cdot d \vec{s}+\mu_0\epsilon_0\dfrac{d }{d t}\int_S \vec{E}\cdot d \vec{s}
\end{equation}

y a su vez

\begin{gather}
    \vec{\nabla}\cdot \vec{E}= \dfrac{\rho}{\epsilon_0}\\
    \vec{\nabla}\cdot \vec{B}= 0\\
    \vec{\nabla}\times\vec{E}=-\dfrac{\partial\vec{B}}{\partial t}\\
    c^2\vec{\nabla}\times\vec{B} =\dfrac{\vec{J}}{\epsilon_0}+\dfrac{\partial
\vec{E}}{\partial t}\label{eq:maxwell3}\\
    \vec{\nabla}\times \vec{B}=\mu_0\vec{J}+\mu_0 \epsilon_0 \dfrac{\partial \vec{E}}{\partial t}
\end{gather}